\chapter{Objetivos y Alcance}

Este capítulo concreta los objetivos que guían el desarrollo del proyecto y delimita el alcance del sistema resultante.
Se estructuran las metas generales y específicas bajo criterios verificables, detallando asimismo las
limitaciones técnicas y las exclusiones explícitas que enmarcan la evaluación del sistema.

\section{Objetivo general}

Diseñar, desarrollar y desplegar una plataforma software \textit{cloud-native}
que integre la consulta y comparación de precios de combustibles y puntos de recarga
eléctrica en España. El sistema incorporará capacidades de análisis histórico, predicción de
precios a corto plazo, recomendación de repostaje basada en rutas y asistencia conversacional,
unificando fuentes de datos heterogéneas en una interfaz de soporte a la decisión accesible para el usuario.

\section{Objetivos específicos y medibles}
Los objetivos específicos desglosan el objetivo general en hitos verificables y acotados,
siguiendo el criterio SMART (\textit{Specific, Measurable, Achievable, Relevant, Time-bound}) \cite{Doran1981}, organizados en torno a cinco ejes funcionales:

\subsection{Ingesta y gestión de datos}

\begin{itemize}
      \item \textbf{OE-01.} Implementar un pipeline de sincronización automática
            con la API REST del Ministerio de Industria (MinITUR) que garantice la
            disponibilidad de datos de precios de combustible fósiles actualizados con una frecuencia mínima
            diaria, asegurando su persistencia histórica en un almacén relacional.

      \item \textbf{OE-02.} Integrar los registros de puntos de recarga eléctrica bajo
            un modelo de datos unificado que permita ejecutar consultas combinadas
            junto a los combustibles convencionales.

      \item \textbf{OE-03.} Diseñar e implementar un flujo de exportación de
            series históricas en formato Apache Parquet hacia Google Cloud Storage
            para servir como base de entrenamiento reproducible de los modelos de aprendizaje automático.
\end{itemize}

\subsection{Consulta y visualización geoespacial}

\begin{itemize}
      \item \textbf{OE-04.} Desarrollar una interfaz web progresiva (PWA)
            con mapa interactivo que permita la búsqueda y filtrado de estaciones
            por criterios geográficos (provincia y municipio) y comerciales (marca y tipo de energía),
            con soporte de internacionalización en castellano, inglés y euskera.

      \item \textbf{OE-05.} Implementar la representación del histórico de
            precios por estación mediante gráficos temporales interactivos.
\end{itemize}

\subsection{Predicción de precios}

\begin{itemize}
      \item \textbf{OE-06.} Desarrollar e integrar un servicio de predicción de precios que
            incorpore el precio del crudo Brent como variable exógena \cite{LAHMIRI2024105008},
            con capacidad de generar predicciones individualizadas por estación para un horizonte temporal de 7 días.
      \item \textbf{OE-07.} Evaluar el rendimiento del algoritmo seleccionado frente a
            alternativas de regresión y modelos base (\textit{baselines}) mediante métricas
            estándar de evaluación ($MAE$, $RMSE$ y $R^2$) \cite{s24124013}.
\end{itemize}

\subsection{Recomendación basada en rutas}

\begin{itemize}
      \item \textbf{OE-08.} Desarrollar un motor de recomendación para trayectos
            entre dos puntos que identifique las estaciones óptimas dentro de un corredor
            geométrico, aplicando una función de puntuación ponderada basada en el coste
            económico y el tiempo de desvío.
      \item \textbf{OE-09.} Garantizar la resiliencia del módulo de enrutamiento
            en el entorno de producción mediante el diseño de políticas de control de
            concurrencia y mecanismos de contingencia (\textit{fallback}) ante la indisponibilidad
            de las APIs geoespaciales externas.
\end{itemize}

\subsection{Arquitectura, despliegue y calidad}

\begin{itemize}
      \item \textbf{OE-10.} Diseñar una arquitectura orientada a microservicios
            desacoplados mediante un API Gateway centralizado, asegurando que cada componente
            sea gobernable, desplegable y escalable de forma independiente \cite{Gartner2026Trends}.

      \item \textbf{OE-11.} Desplegar la plataforma en Google Cloud Run mediante
            la automatización de flujos de integración y entrega continua (CI/CD)
            distribuidos en siete pipelines independientes en Google Cloud Build.

      \item \textbf{OE-12.} Implementar un módulo de identidad basado en el estándar 
      JWT transmitido mediante \textit{cookies} seguras (\textit{httpOnly}), compatible 
      con autenticación federada a través de Google OAuth 2.0.
      
      \item \textbf{OE-13.} Integrar un asistente de voz basado en un pipeline síncrono 
      (STT $\rightarrow$ LLM $\rightarrow$ TTS) utilizando la API de Gemini \cite{gemini_2_5}, 
      habilitando capacidades de \textit{tool calling} \cite{genai_tool_calling} para interactuar 
      con los servicios internos a través del gateway.
\end{itemize}

\section{Alcance del sistema}
\label{sec:alcance}

El alcance define los límites funcionales y técnicos de la solución entregada, delimitando las capacidades operativas del sistema.

\subsection{Ámbito geográfico y de datos}
La plataforma cubre la totalidad del territorio español. 
La ingesta de datos de carburantes utiliza la API pública del Ministerio de Industria, Comercio y Turismo 
\cite{datosgob_carburantes}, que publica precios de aproximadamente 11.000 estaciones de servicio activas
en España. Los puntos de recarga eléctrica se obtienen del portal
mapareve.es, complementando la oferta de datos de movilidad eléctrica
disponible institucionalmente \cite{reve_mapa}.

\subsection{Componentes funcionales incluidos}

La plataforma especificada comprende los siguientes componentes funcionales:

\begin{enumerate}
      \item \textbf{Servicio de gasolineras:} ingesta, normalización,
            persistencia y consulta de precios de carburantes con historial
            temporal. Ofrece búsqueda geoespacial por radio, filtrado por
            provincia, municipio, marca y tipo de combustible, e integra
            los puntos de recarga eléctrica en el mismo modelo de datos.

      \item \textbf{Servicio de usuarios:} registro, autenticación
            (tradicional y OAuth Google), gestión de perfiles con
            preferencias de combustible y vehículo, y sistema de favoritos
            sincronizado. Emite credenciales mediante JWT almacenado en
            \textit{httpOnly cookie} y expone endpoints internos para la
            comunicación entre servicios.

      \item \textbf{Servicio de recomendación:} cálculo de las gasolineras
            óptimas para un trayecto A$\rightarrow$B definido por el usuario,
            mediante identificación de candidatas en el corredor geográfico
            de la ruta, cálculo de desvíos reales respecto al trayecto
            original, evaluación mediante función de puntuación ponderada
            por precio y desvío, y presentación de un ranking con estimación
            del ahorro económico. Soporta configuración del tipo de
            combustible, desvío máximo aceptable y capacidad del depósito,
            e integra un backend de enrutamiento externo para el cálculo
            preciso de distancias y tiempos de desvío.

      \item \textbf{Servicio de predicción:} generación de predicciones de
            precio a corto plazo por estación, basadas en modelos de
            \textit{gradient boosting} entrenados sobre series históricas
            de precios. Incorpora el precio del crudo Brent como variable
            exógena y cubre las estaciones con historial suficiente y
            mayor demanda de uso.

      \item \textbf{Asistente de voz:} interfaz conversacional en lenguaje
            natural que permite al usuario consultar precios y estaciones
            cercanas mediante voz. Implementa un pipeline
            STT$\rightarrow$LLM$\rightarrow$TTS con capacidad de
            \textit{tool calling} para acceder en tiempo real a los datos
            del sistema a través del API Gateway.

      \item \textbf{API Gateway:} punto de entrada único para el frontend,
            que actúa como proxy reverso hacia los servicios internos,
            agrega la documentación OpenAPI de todos ellos, gestiona la
            autenticación mediante cookies y sirve de bridge WebSocket
            para el asistente de voz.

      \item \textbf{Frontend PWA:} aplicación web progresiva instalable
            con mapa interactivo, historial de precios, gestión de
            favoritos, recomendación de ruta y asistente de voz integrado.
            Ofrece soporte \textit{offline}, modo oscuro e
            internacionalización en tres idiomas
            \cite{web_dev_pwa,i18next}.
\end{enumerate}


\subsection{Entorno de despliegue}

La plataforma se aloja en la nube. El trabajo exploratorio ha permitido seleccionar Google Cloud Platform (GCP) \cite{cloud_run2026} como entorno principal,
aprovechando su oferta de servicios serverless y gestionados que facilitan el desarrollo y escalado de aplicaciones cloud-native. La arquitectura se basa en
microservicios alojados en Cloud Run para la ejecución de contenedores, Artifact Registry \cite{artifact_registry} para el almacenamiento de imágenes
Docker \cite{docker}, Cloud Build \cite{cloud_build2026} para la automatización del pipeline CI/CD y Secret Manager \cite{secret_manager} para la gestión
segura de credenciales. Por último, la base de datos relacional se delega en Neon (PostgreSQL gestionado), mientras que el almacenamiento de objetos para
el pipeline de ML se realiza en Google Cloud Storage \cite{cloud_storage2026}.

\section{Limitaciones}

Las limitaciones recogen aquellos aspectos que, estando dentro del alcance
conceptual del proyecto, presentan restricciones técnicas o de recursos
conocidas en la versión entregada.

\begin{itemize}
      \item \textbf{Frecuencia de actualización de precios:} La API del
            Ministerio de Industria publica datos con una frecuencia que puede
            variar entre actualizaciones horarias y diarias según la estación,
            por lo que el sistema no puede garantizar precios en tiempo real
            estricto, sino \textit{near real-time} con un desfase máximo
            determinado por la fuente oficial.
      \item \textbf{Cobertura del modelo predictivo:} La generación y mantenimiento de modelos predictivos por estación es intensiva
            en recursos computacionales (entrenamiento, almacenamiento de artefactos y ejecución de inferencia periódica) y
            depende de la disponibilidad de series temporales suficientemente largas y de calidad. Por ello, el servicio de
            predicción limita su cobertura a las estaciones que disponen de historial suficiente y a las marcadas como favoritas por los usuarios,
            priorizando la relación coste/beneficio y la utilidad práctica frente a una cobertura completa que sería costosa e innecesaria para
            estaciones con datos escasos. El diseño detallado del pipeline de entrenamiento, la política de refresco y las estimaciones de coste
            computacional se documentan en el Capítulo 8 (Sección \ref{subsec:prediccion_impl}).

      \item \textbf{Precisión de la recomendación por ruta:} La calidad
            de la recomendación depende de la disponibilidad y precisión del
            backend de enrutamiento seleccionado. El fallback a distancia
            en línea recta, activable mediante configuración, reduce la
            precisión del corredor geométrico calculado.

      \item \textbf{Escalabilidad del asistente de voz:} El servicio de
            voz está sujeto a los límites de cuota de la API de Gemini y al
            rate limiting configurado (40 peticiones por minuto por IP),
            lo que limita su uso concurrente intensivo.

      \item \textbf{Disponibilidad de datos de recarga eléctrica:}
            La información de puntos de recarga eléctrica procede de una
            fuente de terceros (mapareve.es) cuya disponibilidad y
            completitud no están garantizadas de forma institucional,
            a diferencia de los datos de carburantes del Ministerio.
\end{itemize}

\section{Exclusiones}

Las exclusiones delimitan explícitamente lo que queda fuera del alcance
del presente proyecto, evitando ambigüedades en la evaluación del
trabajo realizado.

\begin{itemize}
      \item \textbf{Aplicación móvil nativa:} El sistema se entrega como
            PWA instalable, pero no se desarrolla una aplicación nativa para
            iOS o Android. La PWA cubre los casos de uso de movilidad mediante
            el navegador del dispositivo.

      \item \textbf{Integración con navegadores GPS:} No se integra con
            sistemas de navegación en tiempo real tipo Google Maps o Waze para
            guía turn-by-turn. La recomendación de ruta se limita al cálculo
            de la estación óptima, no a la navegación posterior. Sin embargo, la aplicacion
            proporciona enlaces directos a Google Maps para la navegación hacia la estación recomendada.

      \item \textbf{Pagos y reservas:} El sistema no contempla
            funcionalidades de pago, reserva de surtidores ni integración
            con sistemas de fidelización de las cadenas de gasolineras.

      \item \textbf{Predicción de precios de electricidad:} El módulo
            predictivo se centra exclusivamente en combustibles líquidos
            (gasolina y diésel). La predicción del coste de recarga eléctrica,
            que depende de tarifas horarias (PVPC) y contratos de operador,
            queda fuera del alcance de este Proyecto de Fin de Grado.

      \item \textbf{Soporte multipaís:} La plataforma está diseñada
            y configurada para el mercado español. La extensión a otros mercados
            europeos requeriría adaptaciones en las fuentes de datos,
            la normalización de formatos y la cobertura geográfica de los
            servicios de enrutamiento.

      \item \textbf{Panel de administración:} No se desarrolla una
            interfaz de administración gráfica para la gestión de usuarios,
            sincronizaciones o modelos ML. La administración se realiza
            mediante endpoints protegidos por \textit{internal secret}
            y acceso directo a la base de datos.
\end{itemize}